% Generated draft: requires mathematical and source review.
\documentclass[11pt]{article}
\usepackage[margin=25mm]{geometry}
\usepackage{fontspec}
\setmainfont{Latin Modern Roman}
\usepackage{amsmath,amssymb,mathtools,unicode-math}
\setmathfont{Latin Modern Math}
\usepackage{enumitem,needspace,xurl,hyperref}
\setlength{\parindent}{0pt}
\setlength{\parskip}{5pt}
\setlength{\emergencystretch}{3em}
\newcommand{\N}{\mathbb{N}}
\newcommand{\Z}{\mathbb{Z}}
\newcommand{\Q}{\mathbb{Q}}
\newcommand{\R}{\mathbb{R}}
\newcommand{\C}{\mathbb{C}}
\newcommand{\F}{\mathbb{F}}
\newcommand{\A}{\mathbb{A}}
\newcommand{\PP}{\mathbb{P}}
\newcommand{\E}{\mathbb{E}}
\newcommand{\eps}{\varepsilon}
\newcommand{\dd}{\,\mathrm{d}}
\begin{document}
\section*{UnsolvedMath ID: 701336}
\textbf{Source problem code:} catalogue entry
\subsection{Definitions and mathematical statement}
For a positive integer nilpotency-class bound \(c\), let \(N_c\) denote the class of torsion-free nilpotent groups of class at most \(c\), called a quasivariety in the source. Explicitly, torsion-freeness means that \(x^n=1\), for \(n\geq 1\), implies \(x=1\). Define the lower central series by \(\gamma_1(G)=G\) and \(\gamma_{i+1}(G)=[\gamma_i(G),G]\), where \([a,b]=a^{-1}b^{-1}ab\); the class bound means \(\gamma_{c+1}(G)=\{1\}\). A subgroup \(D\leq G\) is divisible if
\[
\forall d\in D\;\forall n\in\Z_{\geq1}\;\exists e\in D:\quad e^n=d.
\]
Using the standard relative dominion definition, set
\[
\operatorname{Dom}^{N_c}_G(D)=
\left\{g\in G:\ \forall K\in N_c\;\forall\text{ group homomorphisms }f,h:G\to K,\quad f|_D=h|_D\Longrightarrow f(g)=h(g)\right\}.
\]
The question attributed to A. I. Budkin is whether, for every such \(c\), every \(G\in N_c\), and every divisible subgroup \(D\leq G\),
\[
\operatorname{Dom}^{N_c}_G(D)=D.
\]
Equivalently, must every \(g\in G\setminus D\) admit a group \(K\in N_c\) and homomorphisms \(f,h:G\to K\) agreeing on \(D\) but satisfying \(f(g)\neq h(g)\)?

The source leaves the integer range of \(c\) implicit; the formulation above uses positive integers. If class zero is permitted, its additional case concerns only the trivial group. The source refers to Problem 16.20 for the dominion definition, but that definition is not included in the supplied record, so the standard definition used here requires comparison with that reference. The supplied registry labels this problem open and separately records an explicitly unverified AI claim of an affirmative solution; neither label establishes a mathematical conclusion.
\subsection{Short English statement}
Must every divisible subgroup of a torsion-free nilpotent group equal its dominion within the same nilpotency-class bound? In other words, can every element outside the subgroup be distinguished by two homomorphisms into another group in that class that agree on the subgroup? The supplied affirmative solution claim is explicitly unverified.
\subsection{Sources}
\begin{itemize}
\item Supplied local UnsolvedMath registry record: permanent ID \(701336\), exact source problem code \texttt{catalogue entry}, title \emph{Kourovka Notebook Problem 17.34}; upstream revision \texttt{d371ee917ac9647fbb06cd6fc529deec6f3dda82}.
\item Attributed to A. I. Budkin, 2010. The supplied record cites \emph{The Kourovka Notebook (Unsolved Problems in Group Theory)}, 17th issue (2010), Problem 17.34, and the 21st issue edition (2026), \texttt{arXiv:1401.0300v48}. The collection description names E. I. Khukhro and V. D. Mazurov as editors. The statement refers to Problem 16.20 for the definition of dominion; that referenced text was not supplied.
\item The supplied discussion reports an explicitly unverified AI claim of an affirmative solution attributed to Achyuth Jayadevan, \emph{Preprint}, 2026: \url{https://zenodo.org/records/22752789}. The supplied review notes say that the October 2026 edition of the 21st issue marks this claim with \(\#\).
\item Supplied status review dated \texttt{2026-10-16}: canonical status \texttt{open}, evidence category \texttt{unrefereed\_claim}. This is supplied dated curation, not an independently verified literature assessment. No cited source or claimed solution has been independently consulted for this draft.
\end{itemize}
\end{document}
